\section{AHSME 1962}
        \begin{enumerate}
        	\item (\textit{AHSME 1962}) The expression $\frac{1^{4y-1}}{5^{-1}+3^{-1}}$ is equal to:


 $\textbf{(A)}\ \frac{4y-1}{8} \qquad  \textbf{(B)}\ 8 \qquad  \textbf{(C)}\ \frac{15}{2} \qquad  \textbf{(D)}\ \frac{15}{8} \qquad  \textbf{(E)}\ \frac{1}{8}$


	\item (\textit{AHSME 1962}) The expression $\sqrt{\frac{4}{3}} - \sqrt{\frac{3}{4}}$ is equal to:


 $\textbf{(A)}\ \sqrt{3}/6 \qquad  \textbf{(B)}\ -\sqrt{3}/6 \qquad  \textbf{(C)}\ \sqrt{-3}/6 \qquad  \textbf{(D)}\ 5\sqrt{3}/6 \qquad  \textbf{(E)}\ 1$


	\item (\textit{AHSME 1962}) The first three terms of an arithmetic progression are $x - 1, x + 1, 2x + 3$, in the order shown. The value of $x$ is:


 $\textbf{(A)}\ -2 \qquad  \textbf{(B)}\ 0 \qquad  \textbf{(C)}\ 2 \qquad  \textbf{(D)}\ 4 \qquad  \textbf{(E)}\ \text{undetermined}$


	\item (\textit{AHSME 1962}) If $8^x = 32$, then $x$ equals:


 $\textbf{(A)}\ 4 \qquad  \textbf{(B)}\ \frac{5}{3} \qquad  \textbf{(C)}\ \frac{3}{2} \qquad  \textbf{(D)}\ \frac{3}{5} \qquad  \textbf{(E)}\ \frac{1}{4}$


	\item (\textit{AHSME 1962}) If the radius of a circle is increased by $1$ unit, the ratio of the new circumference to the new diameter is:


 $\textbf{(A)}\ \pi + 2 \qquad  \textbf{(B)}\ \frac{2\pi + 1}{2} \qquad  \textbf{(C)}\ \pi \qquad  \textbf{(D)}\ \frac{2\pi - 1}{2} \qquad  \textbf{(E)}\ \pi - 2$


	\item (\textit{AHSME 1962}) A square and an equilateral triangle have equal perimeters. The area of the triangle is $9\sqrt{3}$ square inches. Expressed in inches the diagonal of the square is:


 $\textbf{(A)}\ 9/2 \qquad  \textbf{(B)}\ 2\sqrt{5} \qquad  \textbf{(C)}\ 4\sqrt{2} \qquad  \textbf{(D)}\ 9\sqrt{2}/2 \qquad  \textbf{(E)}\ \text{none of these}$


	\item (\textit{AHSME 1962}) Let the bisectors of the exterior angles at $B$ and $C$ of triangle $ABC$ meet at $D$. Then the measure in degrees of angle $BDC$ is:


 $\textbf{(A)}\ \frac{1}{2}(90 - A) \qquad  \textbf{(B)}\ 90 - A \qquad  \textbf{(C)}\ \frac{1}{2}(180 - A) \qquad  \textbf{(D)}\ 180 - A \qquad  \textbf{(E)}\ 180 - 2A$


	\item (\textit{AHSME 1962}) Given the set of $n$ numbers, $n > 1$, of which one is $1 - (1/n)$, and all the others are $1$. The arithmetic mean of the $n$ numbers is:


 $\textbf{(A)}\ 1 \qquad  \textbf{(B)}\ n - \frac{1}{n} \qquad  \textbf{(C)}\ n - \frac{1}{n^2} \qquad  \textbf{(D)}\ 1 - \frac{1}{n^2} \qquad  \textbf{(E)}\ 1 - \frac{1}{n} - \frac{1}{n^2}$


	\item (\textit{AHSME 1962}) When $x^9 - x$ is factored as completely as possible into polynomials and monomials with real integral coefficients, the number of factors is:


 $\textbf{(A)}\ \text{more than } 5 \qquad  \textbf{(B)}\ 5 \qquad  \textbf{(C)}\ 4 \qquad  \textbf{(D)}\ 3 \qquad  \textbf{(E)}\ 2$


	\item (\textit{AHSME 1962}) A man drives $150$ miles to the seashore in $3$ hours and $20$ minutes. He returns from the shore to the starting point in $4$ hours and $10$ minutes. Let $r$ be the average rate for the entire trip. Then the average rate for the trip going exceeds $r$, in miles per hour, by:


 $\textbf{(A)}\ 5 \qquad  \textbf{(B)}\ 4\frac{1}{2} \qquad  \textbf{(C)}\ 4 \qquad  \textbf{(D)}\ 2 \qquad  \textbf{(E)}\ 1$


	\item (\textit{AHSME 1962}) The difference between the larger root and the smaller root of

 \[x^2 - px + (p^2 - 1)/4 = 0\]
is


 $\textbf{(A)}\ 0 \qquad  \textbf{(B)}\ 1 \qquad  \textbf{(C)}\ 2 \qquad  \textbf{(D)}\ p \qquad  \textbf{(E)}\ p + 1$


	\item (\textit{AHSME 1962}) When $\left(1 - \frac{1}{a}\right)^6$ is expanded, the sum of the last three coefficients is:


 $\textbf{(A)}\ 22 \qquad  \textbf{(B)}\ 11 \qquad  \textbf{(C)}\ 10 \qquad  \textbf{(D)}\ -10 \qquad  \textbf{(E)}\ -11$


	\item (\textit{AHSME 1962}) $R$ varies directly as $S$ and inversely as $T$. When $R = \frac{4}{3}$ and $T = \frac{9}{14}$, $S = \frac{3}{7}$. Find $S$ when $R = \sqrt{48}$ and $T = \sqrt{75}$.


 $\textbf{(A)}\ 28 \qquad  \textbf{(B)}\ 30 \qquad  \textbf{(C)}\ 40 \qquad  \textbf{(D)}\ 42 \qquad  \textbf{(E)}\ 60$


	\item (\textit{AHSME 1962}) Let $s$ be the limiting sum of the geometric series $4 - \frac{8}{3} + \frac{16}{9} - \cdots$, as the number of terms increases without bound. Then $s$ equals:


 $\textbf{(A)}\ \text{a number between } 0 \text{ and } 1 \qquad  \textbf{(B)}\ 2.4 \qquad  \textbf{(C)}\ 2.5 \qquad  \textbf{(D)}\ 3.6 \qquad  \textbf{(E)}\ 12$


	\item (\textit{AHSME 1962}) Given triangle $ABC$ with base $AB$ fixed in length and position. As the vertex $C$ moves on a straight line, the intersection point of the three medians moves on:


 $\textbf{(A)}\ \text{a circle} \qquad  \textbf{(B)}\ \text{a parabola} \qquad  \textbf{(C)}\ \text{an ellipse} \qquad  \textbf{(D)}\ \text{a straight line} \qquad  \textbf{(E)}\ \text{a curve not listed here}$


	\item (\textit{AHSME 1962}) Given rectangle $R_1$ with one side $2$ inches and area $12$ square inches. Rectangle $R_2$ with diagonal $15$ inches is similar to $R_1$. Expressed in square inches the area of $R_2$ is:


 $\textbf{(A)}\ 9/2 \qquad  \textbf{(B)}\ 36 \qquad  \textbf{(C)}\ 135/2 \qquad  \textbf{(D)}\ 9\sqrt{10} \qquad  \textbf{(E)}\ 27\sqrt{10}/4$


	\item (\textit{AHSME 1962}) If $a = \log_8 225$ and $b = \log_2 15$, then $a$, in terms of $b$, is:


 $\textbf{(A)}\ b/2 \qquad  \textbf{(B)}\ 2b/3 \qquad  \textbf{(C)}\ b \qquad  \textbf{(D)}\ 3b/2 \qquad  \textbf{(E)}\ 2b$


	\item (\textit{AHSME 1962}) A regular dodecagon ($12$ sides) is inscribed in a circle with radius $r$ inches. The area of the dodecagon, in square inches, is:


 $\textbf{(A)}\ 3r^2 \qquad  \textbf{(B)}\ 2r^2 \qquad  \textbf{(C)}\ 3r^2\sqrt{3}/4 \qquad  \textbf{(D)}\ r^2\sqrt{3} \qquad  \textbf{(E)}\ 3r^2\sqrt{3}$


	\item (\textit{AHSME 1962}) If the parabola $y = ax^2 + bx + c$ passes through the points $(-1, 12)$, $(0, 5)$, and $(2, -3)$, the value of $a + b + c$ is:


 $\textbf{(A)}\ -4 \qquad  \textbf{(B)}\ -2 \qquad  \textbf{(C)}\ 0 \qquad  \textbf{(D)}\ 1 \qquad  \textbf{(E)}\ 2$


	\item (\textit{AHSME 1962}) The angles of a pentagon are in arithmetic progression. One of the angles, in degrees, must be:


 $\textbf{(A)}\ 108 \qquad  \textbf{(B)}\ 90 \qquad  \textbf{(C)}\ 72 \qquad  \textbf{(D)}\ 54 \qquad  \textbf{(E)}\ 36$


	\item (\textit{AHSME 1962}) It is given that one root of $2x^2 + rx + s = 0$, with $r$ and $s$ real numbers, is $3 + 2i$ ($i = \sqrt{-1}$). The value of $s$ is:


 $\textbf{(A)}\ \text{undetermined} \qquad  \textbf{(B)}\ 5 \qquad  \textbf{(C)}\ 6 \qquad  \textbf{(D)}\ -13 \qquad  \textbf{(E)}\ 26$


	\item (\textit{AHSME 1962}) The number $121_b$, written in the integral base $b$, is the square of an integer for


 $\textbf{(A)}\ b = 10,\text{ only} \qquad  \textbf{(B)}\ b = 10 \text{ and } b = 5,\text{ only} \qquad  \textbf{(C)}\ 2 \leq b \leq 10 \qquad  \textbf{(D)}\ b > 2 \qquad  \textbf{(E)}\ \text{no value of } b$


	\item (\textit{AHSME 1962}) In triangle $ABC$, $CD$ is the altitude to $AB$, and $AE$ is the altitude to $BC$. If the lengths of $AB$, $CD$, and $AE$ are known, the length of $DB$ is:


 $\textbf{(A)}\ \text{not determined by the information given} \qquad  \textbf{(B)}\ \text{determined only if } A \text{ is an acute angle} \qquad  \textbf{(C)}\ \text{determined only if } B \text{ is an acute angle} \qquad  \textbf{(D)}\ \text{determined only if } ABC \text{ is an acute triangle} \qquad  \textbf{(E)}\ \text{none of these is correct}$


	\item (\textit{AHSME 1962}) Three machines $P$, $Q$, and $R$, working together, can do a job in $x$ hours. When working alone $P$ needs an additional $6$ hours to do the job; $Q$, one additional hour; and $R$, $x$ additional hours. The value of $x$ is:


 $\textbf{(A)}\ \frac{2}{3} \qquad  \textbf{(B)}\ \frac{11}{18} \qquad  \textbf{(C)}\ \frac{3}{2} \qquad  \textbf{(D)}\ 2 \qquad  \textbf{(E)}\ 3$


	\item (\textit{AHSME 1962}) Given square $ABCD$ with side $8$ feet. A circle is drawn through vertices $A$ and $D$ and tangent to side $BC$. The radius of the circle, in feet, is:


 $\textbf{(A)}\ 4 \qquad  \textbf{(B)}\ 4\sqrt{2} \qquad  \textbf{(C)}\ 5 \qquad  \textbf{(D)}\ 5\sqrt{2} \qquad  \textbf{(E)}\ 6$


	\item (\textit{AHSME 1962}) For any real value of $x$ the maximum value of $8x - 3x^2$ is:


 $\textbf{(A)}\ 0 \qquad  \textbf{(B)}\ \frac{8}{3} \qquad  \textbf{(C)}\ 4 \qquad  \textbf{(D)}\ 5 \qquad  \textbf{(E)}\ \frac{16}{3}$


	\item (\textit{AHSME 1962}) Let $a \textcircled{L} b$ represent the operation on two numbers, $a$ and $b$, which selects the larger of the two numbers, with $a \textcircled{L} a = a$. Let $a \textcircled{S} b$ represent the operation which selects the smaller of the two numbers, with $a \textcircled{S} a = a$. Which of the following three rules is (are) correct?


 (1) $a \textcircled{L} b = b \textcircled{L} a$, (2) $a \textcircled{L} (b \textcircled{L} c) = (a \textcircled{L} b) \textcircled{L} c$,
(3) $a \textcircled{S} (b \textcircled{L} c) = (a \textcircled{S} b) \textcircled{L} (a \textcircled{S} c)$.


 $\textbf{(A)}\ (1)\text{ only} \qquad  \textbf{(B)}\ (2)\text{ only} \qquad  \textbf{(C)}\ (1)\text{ and }(2)\text{ only} \qquad  \textbf{(D)}\ (1)\text{ and }(3)\text{ only} \qquad  \textbf{(E)}\ \text{all three}$


	\item (\textit{AHSME 1962}) The set of $x$-values satisfying the equation $x^{\log_{10} x} = x^3/100$ consists of:


 $\textbf{(A)}\ \frac{1}{10},\text{ only} \qquad  \textbf{(B)}\ 10,\text{ only} \qquad  \textbf{(C)}\ 100,\text{ only} \qquad  \textbf{(D)}\ 10 \text{ or } 100,\text{ only} \qquad  \textbf{(E)}\ \text{more than two real numbers.}$


	\item (\textit{AHSME 1962}) Which of the following sets of $x$-values satisfy the inequality

 \[2x^2 + x < 6?\]
$\textbf{(A)}\ -2 < x < \frac{3}{2} \qquad  \textbf{(B)}\ x > \frac{3}{2} \text{ or } x < -2 \qquad  \textbf{(C)}\ x < \frac{3}{2} \qquad  \textbf{(D)}\ \frac{3}{2} < x < 2 \qquad  \textbf{(E)}\ x < -2$


	\item (\textit{AHSME 1962}) Form I


 Consider the statements: (1) $p \wedge q$ (2) $p \wedge \sim q$ (3) $\sim p \wedge q$ (4) $\sim p \wedge \sim q$, where $p$ and $q$ are statements.


 How many of these imply the truth of $\sim(p \wedge q)$?


 Form II


 Consider the statements: (1) $p$ and $q$ are both true (2) $p$ is true and $q$ is false (3) $p$ is false and $q$ is true (4) $p$ is false and $q$ is false.


 How many of these imply the negation of the statement "$p$ and $q$ are both true"?


 $\textbf{(A)}\ 0 \qquad  \textbf{(B)}\ 1 \qquad  \textbf{(C)}\ 2 \qquad  \textbf{(D)}\ 3 \qquad  \textbf{(E)}\ 4$


	\item (\textit{AHSME 1962}) The ratio of the interior angles of two regular polygons is $3:2$. How many such pairs are there?


 $\textbf{(A)}\ 1 \qquad  \textbf{(B)}\ 2 \qquad  \textbf{(C)}\ 3 \qquad  \textbf{(D)}\ 4 \qquad  \textbf{(E)}\ \text{infinitely many}$


	\item (\textit{AHSME 1962}) If $x_{k+1} = x_k + \frac{1}{2}$ for $k = 1, 2, \cdots, n - 1$ and $x_1 = 1$, find

 \[x_1 + x_2 + \cdots + x_n.\]
$\textbf{(A)}\ \frac{n+1}{2} \qquad  \textbf{(B)}\ \frac{n+3}{2} \qquad  \textbf{(C)}\ \frac{n^2-1}{2} \qquad  \textbf{(D)}\ \frac{n^2+n}{4} \qquad  \textbf{(E)}\ \frac{n^2+3n}{4}$


	\item (\textit{AHSME 1962}) The set of $x$-values satisfying the inequality $2 \leq |x - 1| \leq 5$ is:


 $\textbf{(A)}\ -4 \leq x \leq -1 \text{ or } 3 \leq x \leq 6 \qquad  \textbf{(B)}\ 3 \leq x \leq 6 \text{ or } -6 \leq x \leq -3 \qquad  \textbf{(C)}\ x \leq -1 \text{ or } x \geq 3 \qquad  \textbf{(D)}\ -1 \leq x \leq 3 \qquad  \textbf{(E)}\ -4 \leq x \leq 6$


	\item (\textit{AHSME 1962}) For what real values of $K$ does $x = K^2(x - 1)(x - 2)$ have real roots?


 $\textbf{(A)}\ \text{none} \qquad  \textbf{(B)}\ -2 < K < 1 \qquad  \textbf{(C)}\ -2\sqrt{2} < K < 2\sqrt{2} \qquad  \textbf{(D)}\ K > 1 \text{ or } K < -2 \qquad  \textbf{(E)}\ \text{all}$


	\item (\textit{AHSME 1962}) A man on his way to dinner shortly after $6{:}00$ p.m. observes that the hands of his watch form an angle of $110^\circ$. Returning before $7{:}00$ p.m. he notices that again the hands of his watch form an angle of $110^\circ$. The number of minutes that he has been away is:


 $\textbf{(A)}\ 36\frac{2}{3} \qquad  \textbf{(B)}\ 40 \qquad  \textbf{(C)}\ 42 \qquad  \textbf{(D)}\ 42.4 \qquad  \textbf{(E)}\ 45$


	\item (\textit{AHSME 1962}) If both $x$ and $y$ are integers, how many solutions are there of the equation $(x - 8)(x - 10) = 2^y$?


 $\textbf{(A)}\ 0 \qquad  \textbf{(B)}\ 1 \qquad  \textbf{(C)}\ 2 \qquad  \textbf{(D)}\ 3 \qquad  \textbf{(E)}\ \text{more than } 3$


	\item (\textit{AHSME 1962}) $ABCD$ is a square with side of unit length. Points $E$ and $F$ are taken respectively on sides $AB$ and $AD$ so that $AE = AF$ and the quadrilateral $CDFE$ has maximum area. In square units this maximum area is:


 $\textbf{(A)}\ \frac{1}{2} \qquad  \textbf{(B)}\ \frac{9}{16} \qquad  \textbf{(C)}\ \frac{19}{32} \qquad  \textbf{(D)}\ \frac{5}{8} \qquad  \textbf{(E)}\ \frac{2}{3}$


	\item (\textit{AHSME 1962}) The population of Nosuch Junction at one time was a perfect square. Later, with an increase of $100$, the population was one more than a perfect square. Now, with an additional increase of $100$, the population is again a perfect square. The original population is a multiple of:


 $\textbf{(A)}\ 3 \qquad  \textbf{(B)}\ 7 \qquad  \textbf{(C)}\ 9 \qquad  \textbf{(D)}\ 11 \qquad  \textbf{(E)}\ 17$


	\item (\textit{AHSME 1962}) The medians $AN$ and $BP$ of a triangle with unequal sides are, respectively, $3$ inches and $6$ inches long. Its area is $3\sqrt{15}$ square inches. The length of the third median, in inches, is:


 $\textbf{(A)}\ 4 \qquad  \textbf{(B)}\ 3\sqrt{3} \qquad  \textbf{(C)}\ 3\sqrt{6} \qquad  \textbf{(D)}\ 6\sqrt{3} \qquad  \textbf{(E)}\ 6\sqrt{6}$


	\item (\textit{AHSME 1962}) The limiting sum of the infinite series

 \[\frac{1}{10} + \frac{2}{10^2} + \frac{3}{10^3} + \cdots\]
whose $n$-th term is $n/10^n$ is:


 $\textbf{(A)}\ \frac{1}{9} \qquad  \textbf{(B)}\ \frac{10}{81} \qquad  \textbf{(C)}\ \frac{1}{8} \qquad  \textbf{(D)}\ \frac{17}{72} \qquad  \textbf{(E)}\ \text{larger than any finite quantity}$


\end{enumerate}